\begin{abstract}
Property-based testing is a testing approach in which
a programmer specifies properties
(specifications)
that should always be true
of their functions, and tests these properties
on randomly generated values.
In this work,
we propose a new approach to property-based testing
based on describing the desired function behavior via an eDSL, {\footnotesize{\textsc{SymCheck}}}.
{\footnotesize{\textsc{SymCheck}}} allows writing tests
which focus on describing the behavior of a function
via a declarative language.
We believe this offers some advantages over
traditional property based testing:
the language can be used to specify precise preconditions on inputs and sets of inputs,
a sometimes cumbersome task with traditional test generators.
It can also allow for expressing very precise properties that hold
between a functions input and output--- sometimes, specifying \textit{exactly} which input/output pairs are allowed.
{\footnotesize{\textsc{SymCheck}}} tests can  act as a form of specification
for other programmers trying to read and understand the program.
We conduct case studies in which we apply
{\footnotesize{\textsc{SymCheck}}} to real world data structure libraries.
\end{abstract}

